% TODO proof-reading
\subsection{例外}

私たちの\IT{Month}例（\myref{Java_2D_array_month}）を少し書き直してみましょう：

\begin{lstlisting}[caption=IncorrectMonthException.java,style=customjava]
public class IncorrectMonthException extends Exception
{
	private int index;

	public IncorrectMonthException(int index)
	{
		this.index = index;
	} 
	public int getIndex()
	{
		return index;
	}
}
\end{lstlisting}

\begin{lstlisting}[caption=Month2.java,style=customjava]
class Month2
{
	public static String[] months = 
	{
		"January", 
		"February", 
		"March", 
		"April",
		"May",
		"June",
		"July",
		"August",
		"September",
		"October",
		"November",
		"December"
	};

	public static String get_month (int i) throws IncorrectMonthException
	{
		if (i<0 || i>11)
			throw new IncorrectMonthException(i);
		return months[i];
	};

	public static void main (String[] args)
	{
		try
		{
			System.out.println(get_month(100));
		}
		catch(IncorrectMonthException e)
		{
			System.out.println("incorrect month index: "+ e.getIndex());
			e.printStackTrace();
		}
	};
}
\end{lstlisting}

本質的に、\TT{IncorrectMonthException.class}にはオブジェクトコンストラクタと1つのgetterメソッドがあるだけです。

\TT{IncorrectMonthException}クラスは\TT{Exception}から派生しているため、\TT{IncorrectMonthException}コンストラクタは最初に\TT{Exception}クラスのコンストラクタを呼び出し、次に入力された整数値を唯一の\TT{IncorrectMonthException}クラスフィールドに入れます：

\begin{lstlisting}
  public IncorrectMonthException(int);
    flags: ACC_PUBLIC
    Code:
      stack=2, locals=2, args_size=2
         0: aload_0       
         1: invokespecial #1        // Method java/lang/Exception."<init>":()V
         4: aload_0       
         5: iload_1       
         6: putfield      #2        // Field index:I
         9: return        
\end{lstlisting}

\TT{getIndex()}は単なるgetterです。
\TT{IncorrectMonthException}への\IT{reference}が0番目の\ac{LVA}スロット（\IT{this}）に渡され、\TT{aload\_0}がそれを取得し、\TT{getfield}がオブジェクトから整数値をロードし、\TT{ireturn}がそれを返します。

\begin{lstlisting}
  public int getIndex();
    flags: ACC_PUBLIC
    Code:
      stack=1, locals=1, args_size=1
         0: aload_0       
         1: getfield      #2        // Field index:I
         4: ireturn       
\end{lstlisting}

今度は\TT{Month2.class}内の\TT{get\_month()}を見てみましょう：

\begin{lstlisting}[caption=Month2.class]
  public static java.lang.String get_month(int) throws IncorrectMonthException;
    flags: ACC_PUBLIC, ACC_STATIC
    Code:
      stack=3, locals=1, args_size=1
         0: iload_0       
         1: iflt          10
         4: iload_0       
         5: bipush        11
         7: if_icmple     19
        10: new           #2        // class IncorrectMonthException
        13: dup           
        14: iload_0       
        15: invokespecial #3        // Method IncorrectMonthException."<init>":(I)V
        18: athrow        
        19: getstatic     #4        // Field months:[Ljava/lang/String;
        22: iload_0       
        23: aaload        
        24: areturn       
\end{lstlisting}

オフセット1の\TT{iflt}は\IT{if less than}です。

無効なインデックスの場合、オフセット10の\TT{new}命令を使用して新しいオブジェクトが作成されます。

オブジェクトの型は命令のオペランドとして渡されます（これは\TT{IncorrectMonthException}です）。

その後、そのコンストラクタが呼び出され、インデックスが\ac{TOS}経由で渡されます（オフセット15）。

制御フローがオフセット18に到達すると、オブジェクトは既に構築されているため、今度は\TT{athrow}命令が新しく構築されたオブジェクトへの\IT{reference}を取り、適切な例外ハンドラを見つけるよう\ac{JVM}に信号を送ります。

\TT{athrow}命令はここで制御フローを返しないため、オフセット19には例外ビジネスとは関係ない別の\gls{basic block}があり、オフセット7から到達できます。

ハンドラはどのように動作するのでしょうか？

\TT{Month2.class}の\main：

\begin{lstlisting}[caption=Month2.class]
  public static void main(java.lang.String[]);
    flags: ACC_PUBLIC, ACC_STATIC
    Code:
      stack=3, locals=2, args_size=1
         0: getstatic     #5        // Field java/lang/System.out:Ljava/io/PrintStream;
         3: bipush        100
         5: invokestatic  #6        // Method get_month:(I)Ljava/lang/String;
         8: invokevirtual #7        // Method java/io/PrintStream.println:(Ljava/lang/String;)V
        11: goto          47
        14: astore_1      
        15: getstatic     #5        // Field java/lang/System.out:Ljava/io/PrintStream;
        18: new           #8        // class java/lang/StringBuilder
        21: dup           
        22: invokespecial #9        // Method java/lang/StringBuilder."<init>":()V
        25: ldc           #10       // String incorrect month index: 
        27: invokevirtual #11       // Method java/lang/StringBuilder.append:(Ljava/lang/String;)Ljava/lang/StringBuilder;
        30: aload_1       
        31: invokevirtual #12       // Method IncorrectMonthException.getIndex:()I
        34: invokevirtual #13       // Method java/lang/StringBuilder.append:(I)Ljava/lang/StringBuilder;
        37: invokevirtual #14       // Method java/lang/StringBuilder.toString:()Ljava/lang/String;
        40: invokevirtual #7        // Method java/io/PrintStream.println:(Ljava/lang/String;)V
        43: aload_1       
        44: invokevirtual #15       // Method IncorrectMonthException.printStackTrace:()V
        47: return        
      Exception table:
         from    to  target type
             0    11    14   Class IncorrectMonthException
\end{lstlisting}

ここには\TT{Exception table}があり、オフセット0から11（包含）で例外\\
\TT{IncorrectMonthException}が発生する可能性があり、発生した場合、制御フローがオフセット14に渡されることを定義しています。

実際、メインプログラムはオフセット11で終了します。

オフセット14でハンドラが開始します。ここに到達することはできません。この領域への条件付き/無条件ジャンプはありません。

しかし、\ac{JVM}は例外が発生した場合、実行フローをここに転送します。

最初の\TT{astore\_1}（14で）は例外オブジェクトへの入力\IT{reference}を取り、\ac{LVA}スロット1に格納します。

後で、オフセット31で（この例外オブジェクトの）\TT{getIndex()}メソッドが呼び出されます。

現在の例外オブジェクトへの\IT{reference}がその直前（オフセット30）で渡されます。

残りのコードは文字列操作を行うだけです：
最初に\TT{getIndex()}によって返された整数値が\TT{toString()}メソッドによって文字列に変換され、次に「incorrect month index: 」テキスト文字列と連結され（以前に見たように）、その後\TT{println()}と\TT{printStackTrace()}が呼び出されます。

\TT{printStackTrace()}が終了すると、例外が処理され、通常の実行を続けることができます。

オフセット47には\main関数を終了する\TT{return}がありますが、例外が発生しなかったかのように実行される他のコードがある可能性があります。

IDAが例外範囲をどのように表示するかの例です：

\begin{lstlisting}[caption=著者のコンピューターで見つかったランダムな.classファイルから]
    .catch java/io/FileNotFoundException from met001_335 to met001_360\
 using met001_360
    .catch java/io/FileNotFoundException from met001_185 to met001_214\
 using met001_214
    .catch java/io/FileNotFoundException from met001_181 to met001_192\
 using met001_195
    .catch java/io/FileNotFoundException from met001_155 to met001_176\
 using met001_176
    .catch java/io/FileNotFoundException from met001_83 to met001_129 using \
met001_129
    .catch java/io/FileNotFoundException from met001_42 to met001_66 using \
met001_69
    .catch java/io/FileNotFoundException from met001_begin to met001_37\
 using met001_37
\end{lstlisting}